\section{三角関数の逆数}\label{節:相反三角関数}

三角置換では$\frac1{\cos\theta}$や$\frac1{\sin\theta}$が繰り返し現れる.これらを一つの関数として書くと,平方や微分の関係を読み取りやすい.
\begin{definition}{正割・余割・余接}
分母が$0$でないところで
\[
\sec x=\frac1{\cos x},\qquad
\csc x=\frac1{\sin x},\qquad
\cot x=\frac{\cos x}{\sin x}
\]
と定義する.順に正割,余割,余接という.ここでの逆数と,前節の逆関数は異なる操作である.
\end{definition}

$\sin^2x+\cos^2x=1$を$\cos^2x$,\,$\sin^2x$で割ると,それぞれ
\[
1+\tan^2x=\sec^2x,\qquad
1+\cot^2x=\csc^2x
\]
となる.微分も,すでに知っている商の微分から得られる.例えば
\[
(\sec x)'=\frac{\sin x}{\cos^2x}=\sec x\tan x.
\]
同様に計算して
\begin{equation}\label{式:相反三角関数の微分}
(\csc x)'=-\csc x\cot x,\qquad
(\cot x)'=-\csc^2x
\end{equation}
を得る.この記法により,\,$x=\tan\theta$と置いたときの$\dd x$は$\sec^2\theta\dd\theta$,\,$\sqrt{1+x^2}$は$|\sec\theta|$と短く書ける.ただし,記号を短くしても,絶対値を外すための符号の確認は残る.

\begin{example}{三角関数の逆数の積分}
$\int\sec x\dd x$と$\int\csc x\dd x$を求めよ.
\end{example}

$\sec x$の微分は$\sec x\tan x$,\,$\tan x$の微分は$\sec^2x$である.この二つを足すと
\[
(\sec x+\tan x)'=\sec x(\sec x+\tan x)
\]
となり,元の$\sec x$が因子に現れる.また
\[
(\sec x+\tan x)(\sec x-\tan x)=1
\]
だから,\,$\sec x+\tan x$は定義されるところで$0$にならない.従って
\[
\int\sec x\dd x
=\int\frac{(\sec x+\tan x)'}{\sec x+\tan x}\dd x
=\log|\sec x+\tan x|+C.
\]
これは,分母の微分を分子に作る対数型の積分である.

余割では引き算を考えると
\[
(\csc x-\cot x)'=-\csc x\cot x+\csc^2x
=\csc x(\csc x-\cot x).
\]
さらに$(\csc x-\cot x)(\csc x+\cot x)=1$なので,
\[
\int\csc x\dd x=\log|\csc x-\cot x|+C
\]
となる.どちらも,もとの三角関数が定義される一つの区間で使う.特異点をまたぐ不定積分の定数を,一つにそろえる必要はない.


